% Sections won't have numbers
% Make the page area as large as possible
% Set the first level numbering to arabic and the second level
% to alphabetic. The remaining two levels are arabic
%Change to \pagestyle{empty} to suppress numbers on following pages
%\pagestyle{empty}
% Do not reset outermost enumerate counter


\documentclass{article}
%%%%%%%%%%%%%%%%%%%%%%%%%%%%%%%%%%%%%%%%%%%%%%%%%%%%%%%%%%%%%%%%%%%%%%%%%%%%%%%%%%%%%%%%%%%%%%%%%%%%%%%%%%%%%%%%%%%%%%%%%%%%%%%%%%%%%%%%%%%%%%%%%%%%%%%%%%%%%%%%%%%%%%%%%%%%%%%%%%%%%%%%%%%%%%%%%%%%%%%%%%%%%%%%%%%%%%%%%%%%%%%%%%%%%%%%%%%%%%%%%%%%%%%%%%%%
\usepackage{amsmath}
\usepackage{sw20exm2}
\usepackage[compat2]{geometry}

\setcounter{MaxMatrixCols}{10}
%TCIDATA{OutputFilter=LATEX.DLL}
%TCIDATA{Version=5.50.0.2953}
%TCIDATA{<META NAME="SaveForMode" CONTENT="1">}
%TCIDATA{BibliographyScheme=Manual}
%TCIDATA{Created=Tuesday, October 02, 2012 12:12:00}
%TCIDATA{LastRevised=Tuesday, April 23, 2024 12:34:20}
%TCIDATA{<META NAME="ViewSettings" CONTENT="0">}
%TCIDATA{<META NAME="GraphicsSave" CONTENT="32">}
%TCIDATA{<META NAME="DocumentShell" CONTENT="Exams and Syllabi\SW\SW Exam #1 - 8.5x14 Page">}
%TCIDATA{CSTFile=Exam.cst}

\setlength{\topmargin}{-0.75in}
\setlength{\textheight}{12.25in}
\setlength{\oddsidemargin}{0.0in}
\setlength{\evensidemargin}{0.0in}
\setlength{\textwidth}{6.5in}
\def\labelenumi{\arabic{enumi}.}
\def\theenumi{\arabic{enumi}}
\def\labelenumii{(\alph{enumii})}
\def\theenumii{\alph{enumii}}
\def\p@enumii{\theenumi.}
\def\labelenumiii{\arabic{enumiii}.}
\def\theenumiii{\arabic{enumiii}}
\def\p@enumiii{(\theenumi)(\theenumii)}
\def\labelenumiv{\arabic{enumiv}.}
\def\theenumiv{\arabic{enumiv}}
\def\p@enumiv{\p@enumiii.\theenumiii}
\pagestyle{plain}
\setcounter{secnumdepth}{0}
\newtheorem{theorem}{Theorem}
\newtheorem{acknowledgement}[theorem]{Acknowledgement}
\newtheorem{algorithm}[theorem]{Algorithm}
\newtheorem{axiom}[theorem]{Axiom}
\newtheorem{case}[theorem]{Case}
\newtheorem{claim}[theorem]{Claim}
\newtheorem{conclusion}[theorem]{Conclusion}
\newtheorem{condition}[theorem]{Condition}
\newtheorem{conjecture}[theorem]{Conjecture}
\newtheorem{corollary}[theorem]{Corollary}
\newtheorem{criterion}[theorem]{Criterion}
\newtheorem{definition}[theorem]{Definition}
\newtheorem{example}[theorem]{Example}
\newtheorem{exercise}[theorem]{Exercise}
\newtheorem{lemma}[theorem]{Lemma}
\newtheorem{notation}[theorem]{Notation}
\newtheorem{problem}[theorem]{Problem}
\newtheorem{proposition}[theorem]{Proposition}
\newtheorem{remark}[theorem]{Remark}
\newtheorem{solution}[theorem]{Solution}
\newtheorem{summary}[theorem]{Summary}
\newenvironment{proof}[1][Proof]{\noindent\textbf{#1.} }{\ \rule{0.5em}{0.5em}}
\input{tcilatex}
\begin{document}


\begin{equation*}
\FRAME{itbpF}{1.1381in}{0.761in}{0in}{}{}{Figure}{\special{language
"Scientific Word";type "GRAPHIC";maintain-aspect-ratio TRUE;display
"USEDEF";valid_file "T";width 1.1381in;height 0.761in;depth
0in;original-width 1.1044in;original-height 0.729in;cropleft "0";croptop
"1";cropright "1";cropbottom "0";tempfilename
'SCEIF600.wmf';tempfile-properties "XPR";}}
\end{equation*}%
\ \ \ \ \ \ \ \ \ \ \ \ \ \ \ \ \ \ \ \ \ \ \ \ \ \ \ \ \ \ \ \ \ \ \ \ \ \
\ \ \ \ \ \ \ \ \ \ \ \ \ \textbf{Certamen 1 - Electromagnetismo}

\ \ \ \ \ \ \ \ \ \ \ \ \ \ \ \ \ \ \ \ \ \ \ \ \ \ \ \ \ \ \ \ \ \ \ \ \ \
\ \ \ \ \ \ \ \ \ \ \ \ \ \ \ \ \ \ \ \ \ \ 24/04/2024

\noindent\ \ \ \ \ \ \ \ \ \ \ \ \ \ \ \ \ \ \ \ \ \ \ \ \ \ \ \ \ \ \ \ \ \
\ \ \ \ \ \ \ \ \ \ \ Profesor: Arturo Fern\'{a}ndez P\'{e}rez \hspace*{%
\stretch{0.7500}}

\bigskip

\textbf{Instrucciones: }

\begin{itemize}
\item \emph{Responda de manera clara, ordenada y sin borrones. }

\item \emph{No se revisar\'{a}n respuestas poco claras o ilegibles.}

\item \emph{Puntaje Total: 59 puntos. Puntaje m\'{\i}nimo para la aprobaci%
\'{o}n: 30 puntos.}
\end{itemize}

\bigskip 

NOMBRE:\_\_\_\_\_\_\_\_\_\_\_\_\_\_\_\_\_\_\_\_\_\_\_\_\_\_\_\_\_\_\_\_\_\_%
\_\_\_\_\_\_\_\_\_\_\_\_\_\_\_\_\_\_\_\_\_\_\_\_\_\_\_\_\_\_\_\_\_\_\_\_\_\_%
\_\_\_\_\_\_\_\_\_\_\_\_\_\_\_\_\_\_\_\_\_\_\_\_\_\_\_\_\_\_\_\_\_\_\_\_\_\_

\begin{enumerate}
\item (32 puntos) Considere una l\'{\i}nea de carga positiva con densidad
uniforme $\lambda =5,30$ $\left[ \frac{nC}{m}\right] $, ubicada sobre el eje 
$X$, entre $x=0,10$ [m] y $x=0,30$ [m]. Determine el campo el\'{e}ctrico $%
\vec{E}$ inducido por la l\'{\i}nea de carga en el punto $P$, de coordenadas 
$(-0,07,-0,01)$ [m]. Con su resultado, determine la fuerza el\'{e}ctrica
sobre una carga $q=-4,5$ [$\mu C$] ubicada en el punto $P.$\pagebreak 

\item (27 puntos) Tres cargas puntuales se ubican en el plano cartesiano
como muestra la Fig. 1, cuyas coordenadas se dan en cent\'{\i}metros. Si $%
q_{1}=-3$ [$\mu C$], $q_{2}=-4$ [$\mu C$] y $q_{3}=6$ [$\mu C$]$,$ determine
la fuerza resultante $\vec{F}$ sobre la carga $q_{1}$.%
\begin{equation*}
\FRAME{itbpFU}{2.8911in}{3.0208in}{0in}{\Qcb{Figura 1}}{}{Figure}{\special%
{language "Scientific Word";type "GRAPHIC";maintain-aspect-ratio
TRUE;display "USEDEF";valid_file "T";width 2.8911in;height 3.0208in;depth
0in;original-width 8.4267in;original-height 8.8047in;cropleft "0";croptop
"1";cropright "1";cropbottom "0";tempfilename
'SCELVV01.bmp';tempfile-properties "XPR";}}
\end{equation*}
\end{enumerate}

\end{document}
